\section{Roggenschrotbrot}
\noindent\textsf{\small 1 Laib \textperiodcentered\ Dauer: ca. 3,5 Std. \hfill Brot \& Gebäck}
\vspace{1em}

\begin{ingredientsboxfull}
    \noindent
    % --- Linke Spalte: Brühstück ---
    \begin{minipage}[t]{0.45\linewidth}
        % r = amount aligned right, l = unit aligned left, l = ingridient aligned left
        \begin{tabular}{@{} r l l @{}}
            \multicolumn{3}{l}{\textbf{Brühstück:}} \\
            \midrule
            240,65 & g  & Schrot (Stufe 7.5) \\
            721,95 & ml & Wasser (kochend) \\
            48,13  & g  & Salz \\
            & &
        \end{tabular}
    \end{minipage}%
    \hfill
    \begin{minipage}[t]{0.45\linewidth}
        \begin{tabular}{@{} r l l @{}}
            \multicolumn{3}{l}{\textbf{Hauptteig:}} \\
            \midrule
            1347,63 & g  & Schrot (Stufe 7.5) \\
            828,23  & g  & Schrot (Stufe 3) \\
            1754,73 & ml & Wasser \\
            84,23   & g  & Sauerteigstarter \\
        \end{tabular}
    \end{minipage}
\end{ingredientsboxfull}

\vspace{1.5em}
\noindent\textbf{Zubereitung:}
\vspace{0.5em}

\begin{enumerate}[label=\textbf{\arabic*.}, leftmargin=1.8em, itemsep=1em, topsep=0pt]
    \item \textbf{Brühstück:} Zutaten für das Brühstück mit kochendem Wasser übergießen und klumpenfrei verrühren. Folie direkt auf die Oberfläche der Masse legen und für 8--10 Stunden verquellen lassen.
    
    \item \textbf{Sauerteigstarter Füttern 1:} Sauerteigstarter aus dem Kühlschrank nehmen und 1:1:1 füttern und (3--4 Stunden) warten, bis er sich verdoppelt hat. Sauerteigstarter nochmal 1:1:1 füttern und (3--4 Stunden) warten, bis er sich verdoppelt hat.
    
    \item \textbf{Sauerteigstarter Füttern 2:} Sauerteigstarter nochmal 1:1:1 füttern und (3--4 Stunden) warten, bis er sich verdoppelt hat.

     \item \textbf{Formen auskleiden:} Backpapier zurechtschneiden und Formen damit auslegen.
   
    \item \textbf{Hauptteig:} Sauerteigstarter in einer großen Schüssel mit der Wassermenge für den Hauptteig vermischen. Die restlichen Zutaten dazugeben und alles vermischen, bis alles homogen ist, und auf die Backformen aufteilen.
    
    \item \textbf{Gare:} Die Formen abdecken und 24 Stunden bei 18--20\,°C gehen lassen.
    
    \item \textbf{Backen:} wird auf der folgenden Seite fortgesetzt
    \begin{enumerate}
    	\item Backblech auf unterste Schiene in den Ofen einschieben und auf 250\,°C vorheizen.
    	\item Kastenformen mit Teig auf Gitter direkt über dem Backblech einschieben, 100\,ml kochendes Wasser auf das Blech kippen und den Ofen sofort schließen.
    	\item Nach 15 Minuten den Ofen zum Dampfablassen öffnen und das Backblech herausnehmen. Temperatur auf 200\,°C stellen.
    	\item Nach weiteren 65 Minuten ein Stechthermometer in die Brote stecken und weiterbacken, bis die Kerntemperatur mindestens 98\,°C erreicht.
    	\item Brote aus den Formen nehmen und 10 Minuten weiterbacken.
    	\item Ofen ausmachen, die Klappe leicht öffnen und die Brote mindestens 30 Minuten im Ofen abkühlen lassen.
    \end{enumerate}
         
\end{enumerate}

\vspace{2em}

\begin{tcolorbox}[colback=yellow!5, colframe=yellow!40, boxrule=0.5pt, sharp corners]
    \noindent\textbf{Interaktive time plan:} \\
    When isn the bread supposed to be done? \\
    Date (Format \texttt{DD.MM.YYYY}): 
    \Form{\TextField[name=endDate, width=6em, value={24.10.2026}, align=1]{}} 
    \hspace{1em}
    Time (Format \texttt{HH:MM}): 
    \Form{\TextField[name=endTime, width=4em, value={18:00}, align=1]{}}
\end{tcolorbox}

\vspace{1em}
\noindent\textbf{Calculated schedule:}
\vspace{0.5em}

\noindent
\begin{tabular}{r @{\hspace{1em}} c @{\hspace{1em}} p{\dimexpr\linewidth-10.5em\relax}}
    \timelineField{step1}{3610} & \textcolor{gray}{\textbullet} & \textbf{Brühstück:} Brühstück ansetzen und mit Folie abdecken\\
    \vphantom{M} & \textcolor{gray}{|} & \\
    \timelineField{step2}{3490} & \textcolor{gray}{\textbullet} & \textbf{Sauerteig auffrischen:} Sauerteigstarter füttern \\
    \vphantom{M} & \textcolor{gray}{|} & \\
    \timelineField{step3}{3250} & \textcolor{gray}{\textbullet} & \textbf{Sauerteig auffrischen:} Sauerteigstarter nochmal füttern\\
    \vphantom{M} & \textcolor{gray}{|} & \\
    \timelineField{step4}{3030}  & \textcolor{gray}{\textbullet} & \textbf{Backpapier:} Formen mit Backpapier auskleiden \\
    \vphantom{M} & \textcolor{gray}{|} & \\
    \timelineField{step5}{3010}  & \textcolor{gray}{\textbullet} & \textbf{Hauptteig:} Alles zusammen mischen und in Formen füllen \\
    \vphantom{M} & \textcolor{gray}{|} & \\
    \timelineField{step6}{3000}  & \textcolor{gray}{\textbullet} & \textbf{Lange Gare:} Formen abdecken und 24 Std. reifen lassen \\
    \vphantom{M} & \textcolor{gray}{|} & \\
    \timelineField{step7}{1585}   & \textcolor{gray}{\textbullet} & \textbf{Backofen vorheizen:} Ofen mit Blech auf 250\,°C bringen \\
    \vphantom{M} & \textcolor{gray}{|} & \\
    \timelineField{step8}{1560}   & \textcolor{gray}{\textbullet} & \textbf{Einschießen:} Brote mit Dampf backen (Temperatur 200\,°C) \\
    \vphantom{M} & \textcolor{gray}{|} & \\
    \timelineField{step9}{1470}   & \textcolor{gray}{\textbullet} & \textbf{Ofen abschalten:} Ofen ausmachen und leicht "offnen und Brote weitere 30 Minuten im Ofen lassen\\
    \vphantom{M} & \textcolor{gray}{|} & \\
    \timelineField{step10}{1440}    & \textcolor{gray}{\textbullet} & \textbf{Abk"uhlen:} Brote aus dem Ofen nehmen und abkühlen lassen. \\
    \vphantom{M} & \textcolor{gray}{|} & \\
    \timelineField{step11}{0}    & \textcolor{gray}{\textbullet} & \textbf{Fertig!:} Fertig zum Essen. \\
\end{tabular}

